\chapter{Conclusões e Trabalhos Futuros}\label{chap:conclusoes-e-trabalhos-futuros}

Investigou-se se gerar uma representação comprimida do processo, expandida depois por um programa determinístico, produz modelos BPMN 2.0 válidos de forma mais confiável e mais econômica que a geração direta de XML (Seção~\ref{sec:objetivo-geral}). Para isso, foram construídos uma linguagem de domínio específico, a BPMN-DSL, um transpilador dessa linguagem para o XML do BPMN 2.0, um conjunto de pares de texto e DSL e um modelo pequeno especializado nessa notação, comparados com a geração direta de XML por modelos de grande porte em 46 processos do conjunto de teste.

A resposta é dividida. Quanto à confiabilidade, ela é negativa: com os mesmos modelos de grande porte, pedir a DSL no lugar do XML reduziu a validade de 97,8\% para 62,3\% no DeepSeek~V4~Pro e de 93,5\% para 52,9\% no GLM~5.2, e nenhuma configuração com a DSL alcançou a geração direta de XML (Seção~\ref{sec:res-validade}). Quanto à economia, ela é positiva: a DSL deixa de emitir cerca de 81\% dos \textit{tokens} do XML equivalente (Seção~\ref{sec:res-compressao}), e o modelo pequeno especializado tem o menor custo por saída válida entre todos os braços (Seção~\ref{sec:res-custo}).

A queda de validade tem uma causa identificada. A maior parte das falhas estava nas referências a elementos já escritos, que a notação exigia marcar com um identificador antes de o modelo saber que voltaria a eles. Com a referência pelo rótulo do elemento, o mesmo modelo de grande porte passou de 62,3\% para 87,0\% de validade, ainda abaixo do XML direto. Esse resultado é exploratório, porque a correção foi feita depois de vistos os resultados, mas a mesma falha aparece em todos os modelos pequenos ajustados, cujos dados de treino só usam a referência por identificador (Seção~\ref{sec:res-modelo-especializado}).

O modelo pequeno especializado, por fim, não substitui os de grande porte, mas oferece uma troca. Ele produziu saídas válidas na mesma proporção que o método executado pelo modelo de grande porte, 87,0\%, sem a gramática no \textit{prompt}, a um custo por saída válida 5,5 vezes menor que o desse braço e 9,4 vezes menor que o da geração direta de XML. Em troca, a validade fica abaixo da geração direta de XML e, pela medida exploratória, a fidelidade fica cerca de um ponto abaixo, na escala de 0 a 5, da dos modelos de grande porte (Seção~\ref{sec:res-sintese}).

\section{Contribuições}\label{sec:contribuicoes}

\textbf{A BPMN-DSL e o transpilador.} Uma notação compacta para processos BPMN, com gramática formal, e um transpilador determinístico que a expande em XML do BPMN 2.0. Todas as 1.021 amostras do corpus transpiladas foram aceitas pelo esquema oficial, e o XML gerado tem, em média, 5,37 vezes mais \textit{tokens} que a DSL de origem (Seções~\ref{sec:dsl} e~\ref{sec:transpilacao}).

\textbf{Um conjunto de pares de texto e DSL.} São 768 pares de descrição textual e DSL, filtrados pela reconstrução exata do processo extraído, separados em treino e validação por documento de origem, e um modelo de 7 bilhões de parâmetros ajustado sobre eles por US\$~0,47 (Seções~\ref{sec:dados} e~\ref{sec:sft}).

\textbf{Evidência sobre a forma de saída.} A comparação controlada, com o mesmo modelo e as mesmas instruções e mudando só o formato pedido, mostra que modelos de grande porte já escrevem o XML do BPMN com poucos erros, e que a vantagem de uma notação compacta para esses modelos está no custo, e não na validade.

\textbf{Uma condição para notações destinadas a modelos de linguagem.} Uma notação que será gerada da esquerda para a direita não deve exigir que o gerador anuncie, ao escrever um elemento, que ele será referenciado depois. A referência deve poder apontar para o que já está escrito, usando o que está no contexto do gerador, como o rótulo. A condição não depende da BPMN e vale para qualquer notação com referências cruzadas, embora tenha sido testada aqui de forma exploratória e em um único modelo.

\textbf{Uma avaliação que declara os limites das próprias medidas.} As duas medidas de fidelidade usadas foram submetidas a critérios de validação fixados antes de sua aplicação, e nenhuma passou: a comparação com a referência não concordou com as notas de especialista e foi sensível a paráfrases dos rótulos; o juiz também não concordou com essas notas e, embora insensível a paráfrases, não separou com segurança os modelos corretos dos alterados (Seção~\ref{subsec:fidelidade}). Por isso, as conclusões principais se apoiam na validade e no custo, a fidelidade é reportada como exploratória, e os empates são afirmados apenas com margem de equivalência (Seção~\ref{sec:protocolo}). Os desvios do protocolo inicial estão listados com a data e a direção em que poderiam mover os resultados (Seção~\ref{sec:desvios}).

\section{Limitações}\label{sec:limitacoes}

As ameaças à validade estão detalhadas na Seção~\ref{sec:ameacas}. As que mais limitam as conclusões são cinco.

\textbf{Fidelidade sem medida validada.} Nenhuma medida de fidelidade passou nos critérios de validação. As diferenças de fidelidade entre os braços são indicações, e não confirmações, e a escala do juiz não tem tradução conhecida em número de defeitos.

\textbf{Poucos itens de teste.} Com 46 itens, o intervalo das diferenças de validade é mais largo que a margem de equivalência. Por isso, a igualdade de validade entre o modelo pequeno e o método no modelo de grande porte não pôde ser afirmada como equivalência. Parte dos itens, além disso, tem modelos de referência feitos por iniciantes e textos vagos.

\textbf{Dados de treino de uma única fonte.} 94,4\% dos pares de treino vêm do \textit{handbook} da GitLab, e o teste é feito sobre processos de outros domínios. O desenho não separa o efeito dessa mudança de domínio do efeito do tamanho do modelo na fidelidade mais baixa do modelo pequeno.

\textbf{Custo com preço hipotético.} Nenhum provedor consultado oferecia o modelo pequeno com adaptador cobrando por \textit{token}, e o preço usado é uma hipótese. O preço de equilíbrio mostra que a conclusão se mantém até cerca do dobro desse valor, mas não elimina a dependência.

\textbf{Alcance.} O estudo cobre uma linguagem de saída, uma língua, um \textit{prompt} por formato e uma família de modelo pequeno no braço principal. A generalização para outras linguagens XML, que motivou o método, não foi testada.

\section{Trabalhos futuros}\label{sec:trabalhos-futuros}

\textbf{Referência por rótulo no modelo especializado.} É a extensão mais diretamente sustentada pelos resultados. A maior parte das falhas dos modelos pequenos ajustados é de referência a identificadores nunca declarados, a classe de falha que a referência por rótulo reduziu no modelo de grande porte. O corpus pode ser reescrito nessa forma de maneira determinística, sem dado novo, e o modelo, retreinado.

\textbf{Uma medida de fidelidade validada.} A avaliação humana cega de uma amostra de saídas, por mais de um avaliador, permitiria calibrar o juiz e verificar se a diferença de fidelidade do modelo pequeno se mantém. A validação por perturbação também pode ser repetida apenas sobre referências confiáveis, o que corrige o erro de desenho registrado na Seção~\ref{sec:desvios}.

\textbf{Mais itens de teste.} Um conjunto de teste maior estreitaria os intervalos e permitiria decidir se a validade do modelo pequeno é equivalente à do método no modelo de grande porte, o que hoje é inconclusivo.

\textbf{Dados de fontes diferentes.} Mais seções do mesmo \textit{handbook} não elevaram a validade (Seção~\ref{sec:res-modelo-especializado}). Falta saber se dados de outras fontes, que trazem outros estilos de escrita e outros domínios, melhoram a validade e a fidelidade do modelo pequeno.

\textbf{Novos braços experimentais.} O desenho deixa comparações em aberto que outros braços poderiam fechar. A mais direta é ajustar o modelo pequeno para escrever o XML, com os mesmos textos, o que separaria o efeito da notação do efeito do ajuste, comparação que os braços atuais não permitem. Outras seguem a mesma lógica: testar a referência por rótulo também no GLM~5.2, já que ela só foi testada no DeepSeek~V4~Pro; variar o \textit{prompt} de cada formato, que foi um só por formato; e repetir o ajuste com modelos de outros tamanhos, para ver como a validade e a fidelidade mudam com o porte do modelo.

\textbf{Aprendizado por reforço e decodificação restrita.} Como a validade é verificável por programa, ela pode servir de recompensa num ajuste por reforço, desde que combinada a um termo que impeça o modelo de favorecer saídas mínimas e válidas sem relação com o texto. A decodificação restrita à gramática, que garante a sintaxe durante a geração, é uma alternativa à mesma falha e não foi comparada.

\textbf{Outras linguagens de marcação.} Aplicar o mesmo método a outra linguagem XML com esquema formal e serialização extensa verificaria se o ganho de custo e a condição sobre as referências são propriedades gerais da abordagem ou particularidades do BPMN.
